\chapter{past\_event、frames、spread q70：三个朴素过滤器}

\section{本章要解决的困惑}

TFI 给方向，但当前策略不会只因为 TFI 单边就无条件吃单。它还看三个过滤器：

\begin{lstlisting}
past_event_25_bps
frames_since_mid_change
entry_cross_bps <= prior-day q70
\end{lstlisting}

这三个过滤器都很朴素。它们不是复杂预测器，而是回答三个简单问题：

\begin{enumerate}[leftmargin=2em]
  \item 价格是不是已经提前释放了？
  \item mid 是不是停滞了足够久？
  \item 现在 cross 的成本是不是太贵？
\end{enumerate}

\section{past\_event\_25\_bps：刚才价格动了多少}

\code{past_event_25_bps} 来自过去约 25 个 mid 历史点。直觉是：

\[
past25_t = 10000 \times \log\frac{m_t}{m_{t-25}}.
\]

当前 flat 条件是：

\[
|past25_t| \le 0.5 \text{ bps}.
\]

人话：

\begin{quote}
如果主动流来了，但价格刚才还没动很多，可能存在尚未释放的压力。
\end{quote}

注意，这只是很粗的切口。它不是完整 path model，也没有回答“未来最大值什么时候到”。

\section{frames\_since\_mid\_change：mid 停了多久}

\code{frames_since_mid_change} 统计 mid 连续多少个 frame 没变。代码中两个阈值常见：

\begin{lstlisting}
frames_threshold = 31.0
overlay_threshold = 59.80000000000018
\end{lstlisting}

一个小例子：

\begin{longtable}{p{0.18\textwidth}p{0.22\textwidth}p{0.28\textwidth}p{0.22\textwidth}}
\toprule
frame & mid & 是否变化 & frames \\
\midrule
100 & 0.1122 & 刚变化 & 0 \\
101 & 0.1122 & 没变 & 1 \\
102 & 0.1122 & 没变 & 2 \\
... & ... & ... & ... \\
134 & 0.1122 & 没变 & 34 \\
\bottomrule
\end{longtable}

如果此时 TFI 单边、past25 flat，那么 frames 高表示“主动流来了，但 mid 还没动”。这就是旧四象限里 stale-release 的直觉。

\section{spread q70 admission：零手续费也不能乱吃}

最容易误会的是：CCUSDT 零手续费，所以执行成本为零。错。

如果现在盘口是：

\begin{lstlisting}
bid = 0.1120
ask = 0.1124
mid = 0.1122
\end{lstlisting}

做多 entry 打 ask，理论上刚入场就比 mid 贵半个 spread：

\[
entry\_cross\_bps
= 10000 \times \log\frac{ask}{mid}
\approx \frac{spread\_bps}{2}.
\]

当前 \code{EntrySpreadAdmissionGate} 的规则是：

\[
entry\_cross\_bps \le q70_{\text{prior day}}.
\]

也就是说，今天某个 entry 是否允许，要用前一天训练/统计出的 q70 门槛判断，而不是用当天未来分布。

\warningline{如果用当天全日分布来给当天 entry 过滤，就会把未来市场状态偷进策略。}

\section{三个过滤器如何组合}

把 TFI 和过滤器合起来，当前直觉可以写成：

\begin{lstlisting}
direction = sign(TFI)
need abs(TFI) >= 1
need trade_window_count > 0
prefer abs(past_event_25_bps) <= 0.5
prefer frames_since_mid_change high
need entry_cross_bps <= prior-day q70
\end{lstlisting}

更人话的版本：

\begin{quote}
最近主动成交方向很单边，市场价格刚才没有明显跟上，盘口 mid 又停了比较久，而且现在吃单成本没有太夸张。
\end{quote}

\section{为什么这些过滤器不等于预测模型}

它们只做 entry 时刻的条件判断，没有建模后续路径分布。它们没有回答：

\begin{itemize}[leftmargin=2em]
  \item entry 后最大 favorable move 是多少；
  \item 最大值什么时候出现；
  \item 峰值后什么时候回到 0bps；
  \item fixed60 是不是合理退出；
  \item Delta/Energy 是否应该改变 sizing。
\end{itemize}

所以不要把当前策略说成完整 path prediction。它只是用几个朴素过滤器挑出一类 stale-release 状态。

\section{手动检查点}

拿一行 frame，按下面顺序判断：

\begin{lstlisting}
1. abs(trade_flow_imbalance) >= 1 ?
2. trade_window_count > 0 ?
3. abs(past_event_25_bps) <= 0.5 ?
4. frames_since_mid_change >= 25 or 31 ?
5. spread_bps / 2 <= prior_day_q70_entry_cross ?
\end{lstlisting}

前四步决定是否可能产生 shadow signal；第五步是 admission gate，决定这个 signal 是否允许进入 capacity ledger。

